% strict-tierdiff: the "2tier以上でnative一括伸縮と差が出る例"
% case -- a line whose stretch demand is filled by the
% TIER ORDER (explicit tier 50 first) instead of the native one-shot
% proportional stretch. Assertions via log probes (runner greps).
\documentclass[pdflatex]{jlreq}
\begin{document}
\makeatletter
\def\ptj@par@verifywd{1}%
% hsize chosen so line 1 must stretch ~23pt. The line contains:
%   - a JFM 」pair glue   (tier -10 stretch)
%   - kanjiskip boundaries (tier 0 stretch, ~9 x 2.5pt)
%   - one explicit \hspace{0pt plus 30pt} (tier 50 stretch)
% Native one-shot proportional stretch would give the hspace
%   23pt * 30/52.5 = ~13.1pt and every kanjiskip boundary
%   ~1.1pt; the tier engine gives the hspace the whole 23pt
% (within its 30pt capacity) and the lower tiers ZERO.
\hsize=158pt
\begin{ptjparagraph}
あい」うえお。\hspace{0pt plus 30pt}かきくけこ。さしすせそ。たちつてと。
\end{ptjparagraph}
% STRICT-PASS probes: dump the tier-50 explicit glue's delta (GD@1@k
% where the GI spec has 30pt stretch) and assert the kanjiskip
% boundaries' deltas are 0.
\count0=0
\loop
  \advance\count0 1
  \ifnum\count0<40
    \expandafter\ifx\csname ptj@par@GD@1@\the\count0\endcsname\relax\else
      \message{TIERTOP: line=1 gk=\the\count0
        delta=\csname ptj@par@GD@1@\the\count0\endcsname
        GI=\csname ptj@par@GI@1@\the\count0\endcsname}%
    \fi
\repeat
\makeatother
x
\end{document}
